% Editable equation; distinguish the geometric approximation from integer-satoshi consensus rules.
\studiotemplate[
title={Supply},
title-two={of bitcoin},
value-one={1},
caption-one={...of the quantity of},
caption-two={bitcoin issued in},
caption-three={each block.},
value-two={2},
caption-four={At block 6,930,000},
caption-five={(33 halvings), the subsidy},
caption-six={reaches zero satoshis.},
value-three={5},
caption-seven={The total amount of},
caption-eight={Bitcoin available},
caption-nine={It'll be the sum...},
caption-twelve={Halving: 210,000 blocks},
caption-thirteen={210.000 blocks (\textasciitilde{}4 years),},
body-one={Satoshi rounding gives},
caption-fourteen={under 21M BTC.},
value-eight={3},
caption-fifteen={...through a progression},
caption-sixteen={geometrically, where i is},
caption-seventeen={the halving event number.},
value-nine={4},
equation={$\displaystyle 210\,000\sum_{i=0}^{32}\frac{50}{2^{i}}\approx21\,000\,000$}
]{official-layout-72d4d45fa1a8}
\nocite{lesson-official-decks,lesson-bitcoin-faq,lesson-bitcoin-core-subsidy}
